\subsection{Outcome success is not epistemic success}

A 2026 multi-domain evaluation makes the paper's central control problem unusually direct. R\'ios-Garc\'ia et al. evaluate LLM-based scientific agents across more than 25,000 runs and distinguish task outcomes from the epistemic structure of the reasoning process \cite{riosgarcia2026}. Their analysis finds that successful-looking research trajectories can still ignore evidence, fail to revise under refutation and rarely build convergent multi-test support. The result does not show that an external state architecture such as \RAKL{} is sufficient to repair model-level scientific reasoning. It does show that workflow completion and outcome accuracy are inadequate promotion criteria. For the present framework, this is a falsifiable design requirement: evidence use, contradiction handling and belief revision must be represented and checked independently of whether an agent happened to reach a successful endpoint.

The verification boundary is also increasingly occupied by end-to-end systems. ScientistOne maintains a chain of evidence through literature review, solution discovery and paper writing and audits score verification, specification violations, references and method--code alignment \cite{meng2026scientistone}. A contemporaneous survey of autonomous research agents similarly identifies a persistent gap between runnable research workflows and reviewer-verifiable claims \cite{ding2026verificationgap}. Consequently, \RAKL{} does not claim novelty for evidence chains, auditability checklists or claim-to-source traceability in isolation. Its residual systems claim is the controlled external state around those functions: typed authority, contextual compatibility and obstruction, negative history, bounded workspace reconstruction, open-world discovery bookkeeping and fail-closed state transitions.

\subsection{Evidence access can be a harder ceiling than reasoning scaffolds}

A final function-first search surfaced a particularly direct empirical comparator. Wang holds a production scientific-valuation agent's backbone fixed while ablating reasoning scaffolds/public tools and access to a curated proprietary evidence substrate \cite{wang2026evidence}. In that small stratified study, the structured reasoning layer improves calibration and audit discipline, while the much larger gain in factual coverage arrives only when the otherwise inaccessible evidence substrate is supplied. The study is domain-specific and its curated gold set is not an independent ground truth, so it does not establish a universal law of AI science. Its methodological implication is nevertheless important for the present architecture: an agent cannot reason its way to evidence that its acquisition boundary never exposes.

This narrows the role of \RAKL{} further. Evidence governance cannot compensate for an impoverished evidence universe; it can only make the limitation explicit, preserve acquisition provenance, and prevent missing evidence from being converted into unsupported authority. The intended advantage of the framework is therefore not that state discipline removes factual ceilings. It is that the ceiling, the search routes that produced it, and the claim types that remain blocked become first-class research state. A real superiority evaluation must manipulate reasoning architecture and evidence access separately rather than crediting one for the other's contribution.
